%hw-2.tex
%%%%%%%%%%%%%%%%%%%%%%%%%%%%%%%%%%%%%%%%%%%%%%%%%%%%%%%%%%%%%%%%%%%%%%
%%%%%%            Math 403 Homework \#2, Spring 2025            %%%%%%
%%%%%%                                                          %%%%%%
%%%%%%                 Instructor: Ezra Miller                  %%%%%%
%%%%%%%%%%%%%%%%%%%%%%%%%%%%%%%%%%%%%%%%%%%%%%%%%%%%%%%%%%%%%%%%%%%%%%
\documentclass[11pt]{article}

\oddsidemargin=17pt \evensidemargin=17pt
\headheight=9pt     \topmargin=26pt
\textheight=564pt   \textwidth=433.8pt
\parskip=1ex        %voffset=-6ex

\usepackage{amsmath,amssymb,graphicx,color,enumitem}%,setspace,mathrsfs}%,hyperref
\setlist[enumerate]{parsep=1.5ex plus 4pt,topsep=0ex plus 4pt,itemsep=0ex}
  \setlist[itemize]{parsep=1.5ex plus 4pt,topsep=-1ex plus 4pt,itemsep=-1.33ex}

\leftmargini=5.5ex
\leftmarginii=3.5ex

%for \marginpar to fit optimally
%hoffset=-1.02in
\setlength\marginparwidth{2.2in}
\setlength\marginparsep{1mm}
\newcommand\red[1]{\marginpar{\linespread{.85}\sf%
		   \vspace{-1.4ex}\footnotesize{\color{red}#1}}}
\newcommand\score[1]{\marginpar{\vspace{-2ex}\color{blue}{#1/3}}\hspace{-1ex}}
\newcommand\extra[1]{\marginpar{\color{blue}{#1/1}}\hspace{-1ex}}
\newcommand\total[2]{\marginpar{\colorbox{yellow}{\huge #1/#2}}}
\newcommand\collab[1]{\marginpar{\vspace{-11ex}\colorbox{yellow}{#1/3}}\hspace{-1ex}}
\newcommand\magenta[1]{\colorbox{magenta}{\!\!#1\!\!}}
\newcommand\yellow[1]{\colorbox{yellow}{\!\!#1\!\!}}
\newcommand\green[1]{\colorbox{green}{\!\!#1\!\!}}
\newcommand\cyan[1]{\colorbox{cyan}{\!\!#1\!\!}}
\newcommand\rmagenta[2][0ex]{\red{\vspace{#1}\magenta{\phantom{:}}\,: #2}}
\newcommand\ryellow[2][0ex]{\red{\vspace{#1}\yellow{\phantom{:}}\,: #2}}
\newcommand\rgreen[2][0ex]{\red{\vspace{#1}\green{\phantom{:}}\,: #2}}
\newcommand\rcyan[2][0ex]{\red{\vspace{#1}\cyan{\phantom{:}}\,: #2}}

%For separated lists with consecutive numbering
\newcounter{separated}

\newcommand{\excise}[1]{}
\newcommand{\comment}[1]{{$\star$\sf\textbf{#1}$\star$}}

%%%%%%%%%%%%%%%%%%%%%%%%%%%%%%%%%%%%%%%%%%%%%%%%%%%
%                                                 %
% TO ENABLE GRADING, DO NOT ALTER ABOVE THIS LINE %
%                                                 %
%%%%%%%%%%%%%%%%%%%%%%%%%%%%%%%%%%%%%%%%%%%%%%%%%%%

%new math symbols taking no arguments
\newcommand\0{\mathbf{0}}
\newcommand\CC{\mathbb{C}}
\newcommand\FF{\mathbb{F}}
\newcommand\NN{\mathbb{N}}
\newcommand\QQ{\mathbb{Q}}
\newcommand\RR{\mathbb{R}}
\newcommand\ZZ{\mathbb{Z}}
\newcommand\bb{\mathbf{b}}
\newcommand\kk{\Bbbk}
\newcommand\mm{\mathfrak{m}}
\newcommand\xx{\mathbf{x}}
\newcommand\yy{\mathbf{y}}
\newcommand\GL{\mathit{GL}}
\newcommand\FL{{\mathcal F}{\ell}_n}
\newcommand\into{\hookrightarrow}
\newcommand\minus{\smallsetminus}
\newcommand\goesto{\rightsquigarrow}

%redefined math symbols taking no arguments
\newcommand\<{\langle}
\renewcommand\>{\rangle}
\renewcommand\iff{\Leftrightarrow}
\renewcommand\implies{\Rightarrow}
\renewcommand\aa{\mathbf{a}}

%to explain about LaTeX commands:
\newcommand\command[1]{\texttt{$\backslash$#1}}

%new math symbols taking arguments
\newcommand\ol[1]{{\overline{#1}}}

%redefined math symbols taking arguments
\renewcommand\mod[1]{\ (\mathrm{mod}\ #1)}

%roman font math operators
\DeclareMathOperator\im{im}

%for easy 2 x 2 matrices
\newcommand\twobytwo[1]{\left[\begin{array}{@{}cc@{}}#1\end{array}\right]}

%for easy column vectors of size 2
\newcommand\tworow[1]{\left[\begin{array}{@{}c@{}}#1\end{array}\right]}


%%%%%%%%%%%%%%%%%%%%%%%%%%%%%%%%%%%%%%%%%%%%%%%%%%%%%%%%%%%%%%%%%%%%%%
\begin{document}%%%%%%%%%%%%%%%%%%%%%%%%%%%%%%%%%%%%%%%%%%%%%%%%%%%%%%
%%%%%%%%%%%%%%%%%%%%%%%%%%%%%%%%%%%%%%%%%%%%%%%%%%%%%%%%%%%%%%%%%%%%%%


\title{\mbox{}\\[-8ex]Math 403 Homework \#2, Spring 2025\\\normalsize
	Instructor: Ezra Miller\\[-2.5ex]}
\author{Solutions by: ...your name...\\[1ex]
Collaborators: ...list those with whom you worked on this assignment...\\
(1 point for each of up to 3 collaborators who also list you)}
\date{Due: 11:59pm Saturday 8 February 2025}

\maketitle

\vspace{-5ex}\collab{}%

\noindent
\textsc{Reading assignments}\total{}{48}% cover Lectures 7 - 11
\begin{enumerate}[itemsep=-1ex]\setcounter{enumi}{6}
\item% 7.
for Thu.~30 January
  \begin{itemize}
  \item{}%
  [Serge Lang, \emph{Linear Algebra}, Chapter XII, \S1--\S2]
  separating hyperplanes
  \item{}%
  [Serge Lang, \emph{Linear Algebra}, Chapter XII, \S3--\S4]
  support hyperplane, extreme~point$\!\!$
  \end{itemize}

\item% 8.
for Tue.~4 February
  \begin{itemize}
  \item{}%
  [Wikipedia, \emph{Grassmannian}] as much as you're willing; thru \S3
  (\emph{``as a set''}) at least
  \item{}%
  [Wikipedia, \emph{Group action}] get a feel for the concept; no need
  to overdo it
  \end{itemize}

\item% 9.
for Thu.~6 February
  \begin{itemize}
  \item{}%
  [multivariable calculus, any source]: open \& closed sets;
  differentiability; chain rule \mbox{$D(f \!\circ\! g) = Df \!\circ\!
  Dg$}; inverse function theorem: $Df$ invertible $\implies f$ locally
  invertible%
  \item{}%
  [Wikipedia, \emph{Topological space}] get a feel; through \S5.1
  (\emph{Metric spaces}), say
  \item{}%
  [Wikipedia, \emph{Differentiable manifold}] read up through \S2
  (\emph{Definition})
  \end{itemize}

\item%10.
for Tue.~11 February
  \begin{itemize}
  \item{}%
  [Lax, 7.Isometry, p.87--89] isometry, orthogonal group
  \item{}%
  [Lax, 7.Complex Euclidean structure, p.95--96] unitary matrices
% relevant sections: Adjoint (p.84), Isometry (p.87), Orthogonal Group (p.89)
  \item{}%
  [Treil, \S 5.6--\S 5.7] unitary transformations, rigid motions
  \end{itemize}

\item%11.
for Thu.~13 February
  \begin{itemize}
  \item{}%
  [Treil, \S 6.1--\S 6.2] Spectral theorem, normal operators
  \end{itemize}
\end{enumerate}

\noindent
\textsc{Exercises}% cover Lectures 5 - 8

\begin{enumerate}
\item\ \score{}%1.
Find the Jordan form of a $5 \times 5$ matrix~$A$ whose sole
eigenvalue is~$3$, given that $A - 3I$, $(A - 3I)^2$, $(A - 3I)^3$,
and $(A - 3I)^4$ have kernels of dimension $2$, $3$, $4$, and~$5$,
respectively.  Do the same thing if the kernel dimensions are $2$,
$4$, $5$, and~$5$, respectively.

\item\ \score{}%2.
Find a Jordan form and a corresponding Jordan basis for the matrix
$$%$$ $
A =
\left[
\begin{array}{@{}rrrr}
   7\ &\ 1 &  2 &  2 \\
   1\ &\ 4 & -1 & -1 \\
  -2\ &\ 1 &  5 & -1 \\
   1\ &\ 1 &  2 &  8 
\end{array}
\right].
$$%$$ $

\item\ \score{}%3.
Let $V$ be a vector space over~$\RR$ or~$\CC$ of finite dimension.
Let $\varphi: V \into W$ be an injective homomorphism.  Show that if
$\mu$ is a norm on~$W$ then $\mu \circ \varphi$ is a norm on~$V$.

\item\ \score{}%4.
Let $\nu$ be a norm on $V = \RR^n$ or~$\CC^n$ with unit $\nu$-ball
$B_\nu = \{\xx \in V \mid \nu(\xx) \leq 1\}$.  Prove that $B_\nu$ is
closed, bounded, convex, and \emph{equilibrated}: $\xx \in B_\nu$ and
$|\alpha| \leq 1 \implies \alpha\xx \in B_\nu$.  Also show that the
origin lies interior to~$B_\nu$.  Conversely, if $B \subseteq V$ is
closed, bounded, convex, equilibrated subset with $0$ in its interior,
then a norm $\nu_B: V \to \RR$ with unit ball~$B$ can be defined by
$\nu_B(\xx) = \inf\{\alpha > 0 \mid \frac 1\alpha \xx \in B\}$.

\item\ \score{}%5.
The Hahn--Banach theorem says that for any norm~$\nu$ on~$\CC^n$ and
any subspace $V \subseteq \CC^n$, if $\varphi: V \to \CC$ is a linear
functional whose maximum magnitude $|\varphi(\xx)|$ for $\xx$ in the
unit $\nu$-sphere of~$V$ is~$\alpha$, then $\varphi$ can be extended
to a linear functional on~$\CC^n$ whose maximum magnitude on the unit
$\nu$-sphere in~$\CC^n$ is~$\alpha$.  Show that the Hahn--Banach
theorem is equivalent to the statement that separating hyperplanes
exist for points outside of any closed, bounded, convex, equilibrated
set that contains~$0$ in its interior.

\item\ \score{}%6.
Let $V$ be the vector space of real (or complex) absolutely convergent
series.  The function $x \mapsto \Vert x \Vert_p =
\big(\sum_{i=1}^\infty |x_i|^p\big){}^{\frac 1p}$ for $p \geq 1$ is a
norm on~$V$.  (You may assume this, although no complaints if you'd
like to prove it.)  Prove that $\Vert x \Vert_p$ is not topologically
equivalent to $\Vert x \Vert_q$ if $p < q$.

\item\ \score{}%7.
Sketch the convex hull of
$\Bigg\{\tworow{1\\2},\tworow{1\\-1},\tworow{1\\3},\tworow{-1\\1}\Bigg\}$
and
$\Bigg\{\tworow{-1\\2},\tworow{2\\3},\tworow{-1\\-1},\tworow{1\\0}\Bigg\}$.

\item\ \score{}%8.
Fix an operator $\varphi : \RR^n \to \RR^n$.  If $X \subseteq \RR^n$
is convex and $\xx$ is an extreme point of~$X$, must $\varphi(\xx)$ be
an extreme point of~$\varphi(X)$?  What if $\varphi$ is assumed
invertible?

\item\ \score{}%9.
Prove that the intersection of finitely many closed halfspaces
in~$\RR^n$ can have only finitely many extreme points.  Must it have
any extreme points at all?  If it does have extreme points, must the
intersection equal the convex hull of its extreme points?

\item\ \score{}%10.
Fix a column vector $\bb\in\RR^n$ and an $n \times n$ real matrix~$A$.
Show that the set of solutions $\{\xx \in \RR^n \mid A\xx \leq \bb\}$
of the inequality $A\xx \leq \bb$ is convex, where $\aa \leq \bb$ if
$a_i \leq b_i$ for all~$i$.  What kind of convex set is it?  (No need
to provide full justification of this last bit.)

\item\ \score{}%11.
A \emph{cone} is a subset $C \subseteq \RR^n$ closed under nonnegative
scaling: $\xx \in C$ implies $\alpha \xx \in C$ for all $\alpha \in
\RR_{\geq 0}$.  A~ray $\RR_{\geq 0}\xx = \{\alpha\xx \mid \alpha \in
\RR_{\geq 0}\} \subseteq C$ is \emph{extreme} if $\xx \in \yy + C
\subseteq C$ implies $\yy \in \RR_{\geq 0}\xx$.  Let $\nu$ be an outer
support vector for a closed convex cone~$C$, so $\<\nu,\xx\> \leq 0$
for all $\xx \in C$.  Assume $\nu$ attains a unique maximum on~$C$.
% at the origin\/~$\0$.
Prove that the minimum angle between~$\nu$ and a unit vector~$\xx \in
C$ occurs when $\xx$ lies on an extreme~ray~of~$C$.

\item\ \score{}\label{kxn}%12.
How many $k \times n$ matrices of rank~$k$ with entries in a
field~$\FF_q$ of size~$q$ are there?  Hint: how many nonzero row
vectors are there in $\FF_q^n$?  How many vectors are linearly
independent from the first one?  How many are independent from the
first two?

\item\ \score{}\label{free}%13.
Show that the action of~$\GL_k(F)$ on the set of $k \times n$ matrices
of rank~$k$ is \emph{free}: if $g \neq g'$ are invertible $k \times k$
matrices and $A$ is a $k \times n$ matrix of rank~$k$, then $gA \neq
g'A$.  Use this (and the previous exercise) to compute the cardinality
of the set $G_k(\FF_q^n)$ of $k$-planes in~$\FF_q^n$.  Hint: use
\#\ref{kxn} twice, once with $k = n$ and once with $k < n$.

\item\ \score{}%14.
Does the first claim in \#\ref{free} remain true if the rank of~$A$ is
less than~$k$?

\item\ \score{}%15.
Use \#\ref{free} to count the number of Sets in the game of Set.
(Note: a Set is an affine line, not a linear subspace of
dimension~$1$, so this is not quite a special case of \#\ref{free}.
There is more than one way around this; I'll leave it to you to find
at least one.)  Confirm your count by counting a different way: a
line~$L$ is determined by any pair of points~on~$L$.

\end{enumerate}


%%%%%%%%%%%%%%%%%%%%%%%%%%%%%%%%%%%%%%%%%%%%%%%%%%%%%%%%%%%%%%%%%%%%%%
\end{document}%%%%%%%%%%%%%%%%%%%%%%%%%%%%%%%%%%%%%%%%%%%%%%%%%%%%%%%%
%%%%%%%%%%%%%%%%%%%%%%%%%%%%%%%%%%%%%%%%%%%%%%%%%%%%%%%%%%%%%%%%%%%%%%


%%% Various Emacs customizations:
%%% Local Variables:
%%% fill-column: 70
%%% indent-tabs-mode: t
%%% TeX-electric-sub-and-superscript: nil
%%% TeX-brace-indent-level: 0
%%% LaTeX-indent-level: 0
%%% LaTeX-item-indent: 0
%%% TeX-master: t
%%% End: